\begin{figure}[!t]
\centering
\includegraphics[width=.9\columnwidth]{figures/fate_api9e.pdf}
\caption{
\label{fig:api}
\textbf{The \taskName{} task.}
The input $x$ is a text passage produced by a generation model.
Our \emph{Research \& Revision} model %
outputs an attribution report $A$ containing retrieved evidence snippets, along with a revision $y$ whose content can be \emph{attributed} to the evidence in $A$ while \emph{preserving} other properties of $x$ such as style or structure.
}
\end{figure}
\section{Introduction} \label{sec:intro}
Generative language models (LMs) and other text generation models are now the backbone of many AI systems. For example, large language models can perform multi-step reasoning \cite{nye2021show, Wei2022COT}, generate plans \cite{Ahn2022SayCan}, use tools and APIs \cite{shin2021constrained,Thoppilan2022LaMDA}, and answer open-domain questions \cite{petroni2019language,roberts2020much}.

Despite these incredible advances, state-of-the-art LMs still frequently produce biased, misleading, or unsupported content, colloquially called ``hallucinations'' \cite{maynez2020faithfulness,menick2022teaching}. To make LMs more trustworthy, we want to justify each generation by an \emph{attribution report} \cite{Rashkin2021AIS,Bohnet2022AttributedQA} that contains supporting evidence from trusted sources (e.g., encyclopedia or articles) where appropriate.

Most existing LMs, such as those based on sequence-to-sequence architectures, lack a built-in mechanism for attribution. Even \emph{retrieval-augmented models} \cite{gu2020realm,lewis2020rag}, which retrieve relevant documents and then condition on them to generate text, still do not guarantee attribution.
Prior work has shown that retrieval-augmented models generate text that either includes additional information outside the retrieved documents \cite{dziri2022origin}, ignores the documents altogether \cite{krishna2021hurdles}, or even contradicts the documents \cite{longpre2021entity}. In fact, occasionally ignoring the retrievals can make the models more robust to bad retrievals \cite{khandelwal2019generalization}, illustrating that end-task performance and attribution are not always aligned.

Instead of constraining LMs to generate attributed text, we propose a model-agnostic approach to improve the attribution of any existing LM: \emph{\approachfull{}} (\approach{}).
The approach is inspired by works on fact-checking\footnote{In this paper, we generally avoid the term ``fact-checking'' other than to reference relevant literature, because we only address attribution, and attribution does not entail correctness. Even if a claim is attributed to a particular source, it does not guarantee that the source is ``correct'' \cite{menick2022teaching}.} where simple research-and-revise workflows are effective at attributing or correcting unattributed claims made by humans \cite{thorne2018fever,schuster2021get,thorne2021evidence}.
As shown in Figure~\ref{fig:api}, after generating text with the LM, \approach{} does \emph{research} to retrieve relevant evidence, and then \emph{revises} the text to make it consistent with the evidence while preserving qualities like style or structure, enabling the revised text to be seamlessly used in place of the original. \approach{} can be viewed as a retrieval-augmented model where retrieval happens \emph{after} generation rather than before. This allows \approach{} to stand on the shoulders of giant LMs without having to modify them to support attribution.

In our effort to expand the scope of Research \& Revision models to handle the output of arbitrary LMs, we make the following contributions.
First, we formalize the \taskName{} \textbf{task} and propose new \textbf{metrics} that evaluate revision models not just on their ability to produce well-attributed revisions, but also on their ability to otherwise \emph{preserve} original properties of the text.
Second, we use these metrics to \textbf{benchmark} how existing revision models perform on various types of LM outputs such as knowledge-intensive statements, reasoning chains, and dialog responses.
Finally, we find that existing revision models do not always generalize across many tasks (and were not originally intended to), and therefore propose a new research-and-revise \textbf{model} that leverages the power of few-shot prompting in large language models to robustly generalize across domains.